\section{The selected limit in the symmetric projective example}
\label{sec:numa-selected-limit}

The equation of the projective example has many solutions. Nevertheless, its normalized positive-exponent approximations have a unique limit. We prove this by identifying every bounded solution with all the symmetries of the data. The argument is specific to this example.

Choose the cutoff defining $\rho$ even, and let $K$ be the group generated by the block unitary transformations $U(2)\times U(2)$ and the block exchange $\iota$. These transformations preserve $\omega$ and $f\omega_X^3$. Denote by $u_*$ the normalized even solution constructed using the symmetric transition pair in \cref{prop:numa-radial-continuum}.

\begin{theorem}\label{thm:numa-selected-limit}
For the symmetric projective data, $u_*$ is the unique normalized bounded $K$-invariant solution of
\[
 \mathrm{MA}_\omega(u)=f\omega_X^3.
\]
If $u_\lambda$ is the unique solution of
\[
 \mathrm{MA}_\omega(u_\lambda)
      =\mathrm{e}^{\lambda u_\lambda}f\omega_X^3,\qquad\lambda>0,
\]
then
\[
 u_\lambda-\sup_Xu_\lambda\longrightarrow u_*
       \quad\hbox{uniformly as }\lambda\downarrow0.
\]
The selected limit is $C^{1,1}$. This is a selection theorem for these data; it does not assert uniqueness of all normalized solutions of the zero-exponent equation.
\end{theorem}

We give the radial measure calculation needed to treat bounded solutions, rather than assume differentiability of a cluster limit. Use the notation $x,y,z,P,Q,R$ of \cref{sec:degenerate-hermitian-nonuniqueness}. For a scalar $s$ write
\[
 A(r,s)=(x(r)+s)(y(r)-s),\qquad
 \Phi(r,s)=\int_0^s A(r,t)\,\mathrm{d}t.
\]
Partial derivatives of $\Phi$ with respect to $r$ below keep $s$ fixed.

\begin{lemma}\label{lem:numa-radial-flux}
Suppose $u$ is a bounded block-unitary-invariant solution of the zero-exponent equation. On $U$ write $u+\rho=G(r)$. Then $G$ is locally Lipschitz, its derivative $h=G'$ has locally bounded variation, and its one-sided representatives satisfy
\[
 -x(r)\leq h(r\pm)\leq y(r),\qquad \mathrm{d}h+z\,\mathrm{d}r\geq0.
\]
In fact $h$ is continuous. With
\[
 k(r)=\begin{cases}0,&|r|\leq1,\\x(r)y(r)z(r),&|r|>1,\end{cases}
\]
it satisfies the distributional identity on $\mathbb R$
\begin{equation}\label{eq:numa-flux-identity}
 \mathrm{d}\bigl(\Phi(r,h(r))\bigr)
 =\bigl(k(r)-A(r,h(r))z(r)
                    +\partial_r\Phi(r,h(r))\bigr)\,\mathrm{d}r.
\end{equation}
On every open interval where $A(r,h(r))>0$, $h$ is smooth and
\[
 h'=\frac{k(r)}{A(r,h)}-z.
\]
\end{lemma}
\begin{proof}
Block-unitary orbits in $U$ are exactly the level sets of $r$: normalize the total squared norm of $(A,B)$, and map the two vectors of the same respective lengths by their unitary groups. Thus invariance gives a scalar $G$.

On a compact $r$-interval, the radial direction annihilates $P,Q$, while $R$ is strictly positive in that direction. Testing positivity of $\omega_0+\ddc G(r)$ in this direction gives $G''+z\geq0$ as a one-dimensional distribution. One may obtain this first by smooth approximation of the scalar distribution and then by testing with nonnegative functions. Adding a twice integrated primitive of $z$ makes $G$ convex. A finite bounded convex function on an interval is locally Lipschitz, and its derivative is locally of bounded variation. The quasi-plurisubharmonic representative is the corresponding continuous representative. The horizontal directions give $x+h\geq0$, $y-h\geq0$ almost everywhere. Continuity of $x,y$ and the one-sided limits of a function of bounded variation give these bounds for both one-sided representatives. The radial inequality is precisely $\mathrm{d}h+z\,\mathrm{d}r\geq0$; in particular jumps of $h$ can only be upward.

We next justify the formula for the measure, including possible jumps. Identify $U$ smoothly with $\mathbb R\times Y$, where $Y=\{r=0\}$, by real scaling of one block. The forms $P,Q$ are independent of $r$ in this identification. So is $\dc r$: in local coordinates it is the angular differential of the relative complex scale plus fixed base terms. Therefore
\[
 P\wedge Q\wedge R=\mathrm{d}r\wedge\nu,
\]
where $\nu$ is a fixed smooth positive volume form on the compact orbit $Y$, with the orientation chosen by this equality.

For smooth $G$, the rank-one wedge formula gives the radial coefficient of its Monge--Amp\`ere measure, after dividing by $6\nu$, as
\[
 A(r,h)(z+h')\,\mathrm{d}r
 =A(r,h)z\,\mathrm{d}r+\mathrm{d}\bigl(\Phi(r,h)\bigr)
                         -\partial_r\Phi(r,h)\,\mathrm{d}r.
\]
This identity passes to a locally Lipschitz $G$ with the preceding lower Hessian bound. Indeed, convolve $G$ on a slightly larger compact interval. The derivatives $h_\epsilon$ are uniformly bounded and converge to $h$ almost everywhere and in every finite $L^q$ space. Their second derivatives have a uniform lower bound. Consequently the functions $G_\epsilon(r)$ have a common quasi-plurisubharmonic lower Hessian bound on compact subsets of $U$, and converge uniformly to $G(r)$. A fixed smooth local correction and bounded Bedford--Taylor continuity give convergence of their signed products. On the displayed right side, the two ordinary terms converge in $L^1$, and $\Phi(r,h_\epsilon)\to\Phi(r,h)$ in $L^1$, giving distributional convergence of its derivative. This proves the same formula for the bounded product of $G$.

The specified smooth data measure has radial coefficient $k(r)\,\mathrm{d}r$, giving \eqref{eq:numa-flux-identity} initially for the almost-everywhere representative of $h$. If $h$ jumped at $r_0$, the atom of its left side would be
\[
 \Phi(r_0,h(r_0+))-\Phi(r_0,h(r_0-))
 =\int_{h(r_0-)}^{h(r_0+)}(x(r_0)+s)(y(r_0)-s)\,\mathrm{d}s.
\]
The right side has no atom. An upward jump of positive size lies in $[-x(r_0),y(r_0)]$, where the integrand is strictly positive in the interior. Since $x+y>0$ everywhere, such an integral is strictly positive. Thus there are no jumps. A locally bounded-variation function whose jumps all vanish has a continuous representative here: its distributional derivative consists of an absolutely continuous measure and a nonnegative singular measure, so its one-sided representatives coincide at every point. We henceforth use this representative. In particular $G$ is $C^1$ and $G'=h$ everywhere.

All coefficients on the right of \eqref{eq:numa-flux-identity} are now continuous, so $\Phi(r,h(r))$ is $C^1$. Where $\partial_s\Phi=A>0$, the implicit function theorem gives $h\in C^1$ and the stated ordinary differential equation. Smoothness then follows by the smooth coefficients of that equation. The function $k$ is smooth, since the original density and its flat joins are smooth.
\end{proof}

\begin{proof}[of \cref{thm:numa-selected-limit}]
Let $u$ be a normalized bounded $K$-invariant solution. Apply the radial lemma. Because $\rho$ and $u$ are invariant under block exchange, $G$ is even, so $h$ is odd and $h(0)=0$.

On $0<r<1$ we have $x=p>0$ and $y=p+b\geq p>0$. Starting from zero, as long as $-p<h<y$, the equation is $h'=-z_0$. Thus $h(r)=-\int_0^r z_0(s)\,\mathrm{d}s$. This decreasing function cannot meet the positive upper boundary $y$. Its first meeting with the lower boundary $-p$ occurs at the unique $t_*\in(0,1)$ satisfying
\[
 \int_0^{t_*}z_0(s)\,\mathrm{d}s=p(t_*),
\]
whose existence and uniqueness were proved in \cref{prop:numa-radial-continuum}.

After that meeting, $h$ must stay on $-p$ until $r=1$. To prove it, suppose a connected component $(a,b)$ of $\{h+p>0\}$ lies to its right. Its left endpoint has $a\geq t_*$ and $h(a)=-p(a)<0$. Near that endpoint, $h<0<y$; hence on the component the interior equation applies for as long as $h<0$. It forces $h'=-z_0<0$, so $h$ cannot first reach zero. The equation therefore holds throughout the component. But then
\[
 (h+p)'=p'-z_0<0,
\]
contradicting positivity of $h+p$ immediately to the right of its zero at $a$. Thus $h=-p$ on $[t_*,1]$. Oddness determines the entire negative half of the band. The slope on $[-1,1]$ is exactly that of the symmetric explicit solution.

We also identify the slope outside the band; boundedness at the projective ends matters here. On $r>1$, $k=xyz>0$. Neither horizontal coefficient can vanish at a finite point of this interval. For the lower one define
\[
 \Psi(r)=\int_{-x(r)}^{h(r)}A(r,s)\,\mathrm{d}s\geq0.
\]
Equation \eqref{eq:numa-flux-identity} and the smooth lower endpoint show that $\Psi$ is $C^1$. If $h(r_0)=-x(r_0)$, its derivative at that point is $k(r_0)>0$: the $Az$ term vanishes, and the two partial-$r$ terms agree at the common endpoint. This contradicts that the nonnegative $C^1$ function $\Psi$ has an interior minimum zero there. The upper-endpoint integral $\int_{h(r)}^{y(r)}A(r,s)\,\mathrm{d}s$ gives the same contradiction, with derivative $-k(r_0)$, when $h=y$. Thus $-x<h<y$ everywhere on $r>1$, and the smooth ordinary differential equation of the lemma applies.

For $r\geq2$ the background is Fubini--Study:
\[
 x=\sigma,\qquad y=1-\sigma,\qquad z=\sigma',
 \qquad \sigma=\frac{\mathrm{e}^r}{1+\mathrm{e}^r}.
\]
Set $a=\sigma+h\in(0,1)$. Its equation is
\[
 a(1-a)a'=\sigma(1-\sigma)\sigma'>0.
\]
Hence $a$ is increasing and has a limit in $(0,1]$ at $+\infty$. If that limit were less than one, $h=a-\sigma$ would approach a negative constant and its primitive $G$ would be unbounded below, contrary to boundedness of $u$ and $\rho$. The limit is therefore one. Integrating the equation gives
\[
 \frac{a^2}{2}-\frac{a^3}{3}
  =\frac{\sigma^2}{2}-\frac{\sigma^3}{3}+C.
\]
Taking the limit at infinity gives $C=0$. The scalar function $s^2/2-s^3/3$ is strictly increasing on $(0,1)$, so $a=\sigma$ and $h=0$ for $r\geq2$. Ordinary uniqueness propagates this zero solution throughout $r>1$, since the coefficients are smooth and the two horizontal coefficients stay positive there. Oddness gives $h=0$ on $r<-1$ as well.

We have proved that $G'$ agrees everywhere with the symmetric explicit slope. Thus $u$ differs from $u_*$ by a constant on $U$. The two quasi-plurisubharmonic functions then agree up to that constant on $X$ by their canonical representatives. Their supremum normalization removes the constant. This proves uniqueness within the bounded $K$-invariant class.

Finally \cref{prop:symmetric-equation-clusters} gives uniform precompactness of the normalized positive-exponent family. Every cluster limit solves the zero-exponent equation and is $K$-invariant, because the data are $K$-invariant. The scalar is one for these data. The uniqueness just proved identifies every cluster limit with $u_*$. If convergence of the whole family failed, a sequence staying a fixed positive uniform distance from $u_*$ would have a convergent subsequence, contradicting that identification. This proves uniform convergence and the asserted $C^{1,1}$ regularity of the limit.
\end{proof}
